\begin{pastexam}[label={px:ntu-109-ee-ds-24}]{臺大電機所 109 學年度（B）第 24 題}
（複選；每題 4 分，每答錯一個選項倒扣 0.8 分至該題 0 分止）

In this exam: The height of a singleton tree is defined as 0. The order of a node in a tree is
defined as the number of its children. The order of a tree is defined as the maximum order among
its nodes.

Which of the following statements are correct?
\begin{enumerate}[label=(\Alph*)]
\item In a binomial heap, there is at most one binomial tree of order 3.
\item In a binomial heap of n items, the height among all binomial trees is O($\log n$).
\item In a binomial heap of n items, there is O($\log n$) binomial trees.
\item Stack is first-in-first-out.
\item Queue is first-in-first-out.
\end{enumerate}
\end{pastexam}

\begin{pastsoln}
答案：(A)、(B)、(C)、(E)。

\begin{itemize}
\item (A) 對。 二項堆積的定義就是「同一個階數最多一棵」（資結書第~8.5~節），合併時同階的兩棵會
連結成高一階的一棵，像二進位進位。 3 階的二項樹 $B_3$ 最多一棵。 （費氏堆積可以有很多棵同階的樹，
但那不是二項堆積。）
\item (B) 對。 $B_k$ 有 $2^k$ 個節點，高度（數邊）是 $k$。 堆積裡最大的樹 $B_k$ 滿足 $2^k \le n$，
所以每棵樹的高度都 $\le \lg n$。
\item (C) 對。 樹的棵數是 $n$ 的二進位裡 $1$ 的個數，最多 $\lfloor \lg n \rfloor + 1$ 棵。
\item (D) 錯。 堆疊是後進先出（LIFO），最後放進去的最先拿出來（資結書第~6.1~節）。
\item (E) 對。 佇列是先進先出（FIFO，資結書第~6.2~節）。
\end{itemize}
\end{pastsoln}
